\section{Conclusion}
\label{sec:conclusion}

We turned negative transfer in gradient-based meta-learning from an accuracy symptom into an explicit, localized, post-hoc explanation. Adaptation gain $\Delta M$ gives a differentiable, sign-aware answer to \emph{whether} adaptation helped, provably dominated by the backbone's own displacement in the fast-adaptation regime; \fama answers \emph{where}, extending Grad-CAM across an entire adaptation trajectory through an $\mathcal{O}(TP)$ adjoint recursion, without ever touching the meta-training set. Two trajectory-aware faithfulness protocols -- a bidirectional deletion test and cascading-randomization/support-intervention sanity checks -- confirm this explanation is causally tied to $S_i$ ($p<0.001$) and collapses exactly where it should when $\theta_0$'s learned content is destroyed, rather than persisting as an architecture-only artifact. Applied across 600 cross-domain tasks, \fama recovers a concrete, visualizable mechanism -- contextual bias, the backbone falling back on background correlations when the object itself is a weak signal -- behind a phenomenon previously known only by its symptom. We hope this reframing of adaptation as an explainable object, rather than a black-box update rule, is a useful step toward meta-learners whose failures can be diagnosed, not just observed.
